\section{Experimental Methodology}\label{sec:methodology}

\subsection{Synthesis}\label{subsec:meth-synthesis}
Every resource and timing figure in Table~\ref{tab:resources} is produced by \texttt{vitis\_hls} running \texttt{csynth\_design} in batch mode against the part \texttt{xcu55c-fsvh2892-2L-e} with a 3.33\,ns clock target, Vitis HLS 2023.2. Reports are committed alongside their generating \texttt{.tcl} scripts so every number is independently regeneratable. These are HLS estimates, not post-route or on-device measurements, for every kernel and variant \emph{except} one: the \texttt{m\_axi}-fixed Shor kernel was additionally taken through \texttt{v++ -c}, \texttt{v++ -l}, and full Vivado implementation, producing the real post-route reports in Table~\ref{tab:postroute} (Section~\ref{subsec:res-hardware}).

\subsection{Hardware test}\label{subsec:meth-hw}
The placed-and-routed \texttt{m\_axi}-fixed Shor kernel's \texttt{.xclbin} was loaded onto a live Xilinx Alveo U55C (XRT 2.15.225) via \texttt{pyxrt}. For each of the 27 single-qubit self-test cases, the kernel was invoked with the packed error vector as a scalar AXI-Lite argument and its result read back through a 4-byte \texttt{xrt::bo} buffer object bound to the kernel's \texttt{m\_axi} output port (synchronised device-to-host after each invocation), not through the AXI-Lite return register that every previous attempt in this line of work depended on. A single opportunistic latency measurement (Section~\ref{subsec:res-hardware}) timed 10,000 further invocations of the same path end to end with host wall-clock timestamps around argument set, kernel start, wait, buffer sync, and buffer read.

\subsection{Software correctness verification}\label{subsec:meth-correctness}
Each kernel has a Python mirror built directly from its HLS source: the Shor mirror parses its 256-entry decoder table out of the \texttt{.cpp} file programmatically rather than by hand, so it cannot silently drift from what HLS actually synthesises. Three checks are run per kernel: (1) an exhaustive single-qubit-Pauli self-test; (2) for the Shor kernel, a bit-for-bit cross-check against an independently written reference decoder (embedded in a separate host driver, built via an unrelated code path) over the complete 512\,\texttimes\,512 input space; (3) for the Steane kernel, exhaustive weight-1 and weight-2 agreement checks across all three decoder modes. None of this constitutes a hardware correctness proof; it establishes that the software model matches the synthesised source's own logic, which is a necessary but not sufficient condition for correct hardware behaviour.

\subsection{Monte Carlo protocol}\label{subsec:meth-mc}
Two Monte Carlo campaigns are reported. The first is phenomenological: for each code and error model, error vectors are drawn directly over the data qubits (bypassing any circuit-level structure) and decoded using the exact per-code failure lookup table obtained by exhaustively enumerating every possible input in advance, so the only sampled quantity is the random error injection itself. \(10^7\) shots are drawn per grid point, at physical error rates \(p \in \{10^{-3}, 5\times10^{-3}, 10^{-2}, 2\times10^{-2}, 5\times10^{-2}, 0.10, 0.15, 0.18, 0.20\}\), seed \texttt{0xC0DE7}. Two noise models are used: single-Pauli (with probability \(p\), exactly one uniformly random qubit receives a uniformly random Pauli) and IID-depolarising (each qubit independently errors with probability \(p\), uniform over \(\{X,Y,Z\}\) conditioned on an error). The repetition code instead uses an IID bit-flip model, since it has no notion of a phase error. Logical error rates are reported with Wilson score 95\% confidence intervals; where three decoder modes are compared, a two-proportion z-test is reported per grid point rather than an unquantified claim of equivalence.

The second campaign is circuit-level: for the repetition code only, Stim's built-in repetition-code memory-experiment circuit generator produces a real multi-round syndrome-extraction circuit under per-gate depolarising noise, measurement-flip noise, and reset-flip noise, decoded with PyMatching's exact minimum-weight perfect matching against the circuit's own detector error model. This is the only circuit-level, multi-round result in this paper; the Shor and Steane codes have no off-the-shelf circuit-level noise model available and are not covered by this campaign.

\subsection{What this methodology does not cover}\label{subsec:meth-gaps}
Beyond the single hardware test above, no result in this paper is measured on FPGA hardware: the repetition-code and Steane kernels have not been placed and routed or tested, and the Shor kernel's latency measurement is a single opportunistic run, not the rigorous, tail-resolving, C++-hosted measurement (Section~\ref{sec:limitations}) a real-time claim would require. The accompanying repository's blocker log records what is needed to extend hardware coverage to the other kernels and to that rigorous measurement: continued access to a shared accelerator card currently in use by other work on this machine, and further place-and-route time.
